\chapter{Conclusion}

Digital Mohalla successfully demonstrates that peer-to-peer neighbourhood resource sharing can be grounded in mathematical trust modeling, formal event-sourcing invariants, and privacy-preserving architecture rather than opaque star ratings or invasive GPS tracking.

At the FYP-1 mid-assessment milestone, the engineering baseline is fully established:
\begin{itemize}
    \item The design baseline (Part A) achieved complete consistency across UML 2.5 domain models, ER schemas, and formal architectural decisions.
    \item Slice A delivers a fully functional, event-sourced borrowing lifecycle in Python 3.12, FastAPI, and PostgreSQL 16, backed by append-only cryptographic event logging, Argon2/JWT authentication, and an Expo presentation client.
    \item Slice B successfully implements the complex computing core: an empirical Bayesian trust engine combining recency-weighted Beta posteriors, 90\% credible intervals, contextual similarity weighting, bounded $k$-hop indirect graph propagation, and transparent decision policies (Rules R1--R9) integrated with a three-layer lender evidence dashboard.
    \item Concurrency safety is mathematically guaranteed at the storage level via PostgreSQL exclusion constraints, completely preventing overlapping reservations.
\end{itemize}

While Slice C (collusion detection and dispute moderation) remains in active development, the platform already establishes a verifiable, reproducible foundation. Digital Mohalla proves that decentralized community sharing can be made mathematically sound, practically usable, and fundamentally respectful of citizen privacy.
